\section{REFERENCIAS}
\begin{flushleft}
	\setlength{\parindent}{0pt}
	\setlength{\parskip}{12pt}
	
	\hangindent=1.25cm\hangafter=1
	Adida, B. (2008). Helios: Web-based open-audit voting. En \textit{Proceedings of the 17th USENIX Security Symposium} (pp. 335--348). USENIX Association. Recuperado de: \url{https://www.usenix.org/legacy/event/sec08/tech/full_papers/adida/adida.pdf}

	\hangindent=1.25cm\hangafter=1
	Ajtai, M. (1996). Generating hard instances of lattice problems. En \textit{Proceedings of the Twenty-Eighth Annual ACM Symposium on Theory of Computing (STOC '96)} (pp. 99--108). ACM. Recuperado de: \url{https://api.semanticscholar.org/CorpusID:14064562}

	\hangindent=1.25cm\hangafter=1
	Anton, H. (2003). \textit{Introducción al álgebra lineal} (3.ª ed.). Limusa Wiley.
	
	\hangindent=1.25cm\hangafter=1
	Baloglu, S., Bursuc, S., Mauw, S., y Pang, J. (2021). Election verifiability revisited: Automated security proofs and attacks on Helios and Belenios. \textit{34th IEEE Computer Security Foundations Symposium (CSF 2021)}, 1--16. Recuperado de: \url{https://eprint.iacr.org/2020/982}

	\hangindent=1.25cm\hangafter=1
	Barrat, J. (2012). El voto electrónico: Entre la técnica y el derecho. \textit{Revista de Derecho Político}, (84), 211--244. Recuperado de: \url{https://doi.org/10.5944/rdp.84.2012.9103}

	\hangindent=1.25cm\hangafter=1
	Benaloh, J. (1987). \textit{Verifiable secret-ballot elections} (Technical Report YALEU/DCS/TR-561) [Doctoral dissertation, Yale University]. Recuperado de: \url{https://www.cs.yale.edu/publications/techreports/tr561.pdf}

	\hangindent=1.25cm\hangafter=1
	Cachero, C. (2003). \textit{CBD-Web: Un método de desarrollo dirigido por modelos para aplicaciones Web} (Tesis doctoral). Universidad de Alicante. Recuperado de: \url{http://hdl.handle.net/10045/10086}

	\hangindent=1.25cm\hangafter=1
	Centro Criptológico Nacional. (2023). \textit{Guía de seguridad de las TIC (CCN-STIC 893): Criptografía post-cuántica}. Centro Criptológico Nacional. Recuperado de: \url{https://www.ccn-cert.cni.es/guias/guias-series-ccn-stic/800-guia-esquema-nacional-de-seguridad/893-ccn-stic-893-criptografia-post-cuantica/file.html}

	\hangindent=1.25cm\hangafter=1
	Chaum, D. (1981). Untraceable electronic mail, return addresses, and digital pseudonyms. \textit{Communications of the ACM}, 24(2), 84--90. Recuperado de: \url{https://api.semanticscholar.org/CorpusID:207604337}

	\hangindent=1.25cm\hangafter=1
	Churata, H. (2020). \textit{Sistema de votación electrónica basado en blockchain para la transparencia electoral universitaria} (Tesis de licenciatura). Universidad Pública de El Alto, La Paz, Bolivia.

	\hangindent=1.25cm\hangafter=1
	Cortier, V., Gaudry, P., y Lesturgeon, Q. (2019). Belenios: A simple private and verifiable electronic voting system. \textit{Foundations of Security, Protocols, and Secrecy}, 1--22. Recuperado de: \url{https://inria.hal.science/hal-02066930v2/file/main.pdf}

	\hangindent=1.25cm\hangafter=1
	Cramer, R., Gennaro, R., y Schoenmakers, B. (1997). A secure and optimally efficient multi-authority election scheme. En \textit{Advances in Cryptology -- EUROCRYPT '97} (pp. 103--118). Springer. Recuperado de: \url{https://api.semanticscholar.org/CorpusID:1653833}

	\hangindent=1.25cm\hangafter=1
	Diffie, W., y Hellman, M. (1976). New directions in cryptography. \textit{IEEE Transactions on Information Theory}, 22(6), 644--654. Recuperado de: \url{https://ee.stanford.edu/~hellman/publications/24.pdf}

	\hangindent=1.25cm\hangafter=1
	Dolev, D., y Yao, A. C. (1983). On the security of public key protocols. \textit{IEEE Transactions on Information Theory}, 29(2), 198--208. Recuperado de: \url{https://api.semanticscholar.org/CorpusID:206556534}

	\hangindent=1.25cm\hangafter=1
	Fernández, C. E. (2021). \textit{Implementación de un sistema de voto electrónico sustentado en tecnología blockchain} (Tesis de licenciatura). Universidad Mayor de San Andrés, La Paz, Bolivia.

	\hangindent=1.25cm\hangafter=1
	Fujioka, A., Okamoto, T., y Ohta, K. (1992). A practical secret voting scheme for large scale elections. En \textit{Advances in Cryptology -- AUSCRYPT '92} (pp. 244--251). Springer. Recuperado de: \url{https://api.semanticscholar.org/CorpusID:2879555}

	\hangindent=1.25cm\hangafter=1
	Fujisaki, E., y Okamoto, T. (2013). Secure integration of asymmetric and symmetric encryption schemes. \textit{Journal of Cryptology}, 26(1), 80--101. Recuperado de: \url{https://api.semanticscholar.org/CorpusID:253648602}

	\hangindent=1.25cm\hangafter=1
	Fúster, A., Caballero, P., De la Guía, C., y Montoya, F. (2004). \textit{Técnicas criptográficas de protección de datos} (2.ª ed.). Ra-Ma Editorial.

	\hangindent=1.25cm\hangafter=1
	Grossman, S. I., y Flores, J. J. (2012). \textit{Álgebra lineal} (7.ª ed.). McGraw-Hill Interamericana.

	\hangindent=1.25cm\hangafter=1
	Hevner, A. R., March, S. T., Park, J., y Ram, S. (2004). Design science in information systems research. \textit{MIS Quarterly}, 28(1), 75--105. Recuperado de: \url{https://www.jstor.org/stable/25148625}

	\hangindent=1.25cm\hangafter=1
	Koch, N., Knapp, A., Zhang, G., y Baumeister, H. (2008). UML-based web engineering. En \textit{Emerging Methods, Technologies, and Process Management in Software Engineering} (pp. 157--191). IGI Global. Recuperado de: \url{https://api.semanticscholar.org/CorpusID:2966831}

	\hangindent=1.25cm\hangafter=1
	Langlois, A., y Stehlé, D. (2015). Worst-case to average-case reductions for module lattices. \textit{Designs, Codes and Cryptography}, 75(3), 565--599. Recuperado de: \url{https://eprint.iacr.org/2012/090.pdf}

	\hangindent=1.25cm\hangafter=1
	Lucena, M. J. (2010). \textit{Criptografía y seguridad en computadores} (4.ª ed.). Universidad de Jaén. Recuperado de: \url{https://www.marcombo.com/descargas/criptografia.pdf}

	\hangindent=1.25cm\hangafter=1
	Lyubashevsky, V. (2012). Lattice signatures without trapdoors. En \textit{Advances in Cryptology -- EUROCRYPT 2012} (pp. 738--755). Springer. Recuperado de: \url{https://eprint.iacr.org/2011/537.pdf}

	\hangindent=1.25cm\hangafter=1
	Lyubashevsky, V., Peikert, C., y Regev, O. (2010). On ideal lattices and learning with errors over rings. En \textit{Advances in Cryptology -- EUROCRYPT 2010} (pp. 1--23). Springer. Recuperado de: \url{https://eprint.iacr.org/2012/230.pdf}

	\hangindent=1.25cm\hangafter=1
	National Institute of Standards and Technology. (2015). \textit{SHA-3 standard: Permutation-based hash and extendable-output functions} (FIPS PUB 202). U.S. Department of Commerce. Recuperado de: \url{https://nvlpubs.nist.gov/nistpubs/FIPS/NIST.FIPS.202.pdf}

	\hangindent=1.25cm\hangafter=1
	National Institute of Standards and Technology. (2024a). \textit{Module-lattice-based key-encapsulation mechanism standard} (FIPS PUB 203). U.S. Department of Commerce. Recuperado de: \url{https://nvlpubs.nist.gov/nistpubs/FIPS/NIST.FIPS.203.pdf}

	\hangindent=1.25cm\hangafter=1
	National Institute of Standards and Technology. (2024b). \textit{Module-lattice-based digital signature standard} (FIPS PUB 204). U.S. Department of Commerce. Recuperado de: \url{https://nvlpubs.nist.gov/nistpubs/FIPS/NIST.FIPS.204.pdf}

	\hangindent=1.25cm\hangafter=1
	Object Management Group. (2014). \textit{Object Constraint Language (OCL), Version 2.4}. OMG Document formal/2014-02-03. Recuperado de: \url{https://www.omg.org/spec/OCL/2.4/PDF}

	\hangindent=1.25cm\hangafter=1
	Panizo, L. (2007). \textit{Seguridad y verificación formal en protocolos de votación electrónica} (Tesis doctoral). Universidad de León, España. Recuperado de: \url{http://hdl.handle.net/10612/886}

	\hangindent=1.25cm\hangafter=1
	Pastor, O., y Molina, J. C. (2007). \textit{Desarrollo de software dirigido por modelos: El método OO-Method}. Editorial Camino.

	\hangindent=1.25cm\hangafter=1
	Pressman, R. S. (2010). \textit{Ingeniería del software: Un enfoque práctico} (7.ª ed.). McGraw-Hill Interamericana.

	\hangindent=1.25cm\hangafter=1
	Ramos, A., y Gómez, P. (2023). Criptografía post-cuántica en sistemas de votación electrónica: Una revisión sistemática. \textit{Revista Iberoamericana de Ciberseguridad y Tecnologías de la Información}, 5(2), 45--62.

	\hangindent=1.25cm\hangafter=1
	Regev, O. (2005). On lattices, learning with errors, random linear codes, and cryptography. En \textit{Proceedings of the Thirty-Seventh Annual ACM Symposium on Theory of Computing (STOC '05)} (pp. 84--93). ACM. Recuperado de: \url{https://cims.nyu.edu/~regev/papers/qcrypto.pdf}

	\hangindent=1.25cm\hangafter=1
	Reniu, J. M. (2008). ¿Votar por Internet? El voto electrónico y sus implicaciones sociopolíticas. \textit{Revista de Estudios Políticos}, (140), 125--154. Recuperado de: \url{https://recyt.fecyt.es/index.php/REP/article/view/37466}

	\hangindent=1.25cm\hangafter=1
	Rivest, R. L., Shamir, A., y Adleman, L. (1978). A method for obtaining digital signatures and public-key cryptosystems. \textit{Communications of the ACM}, 21(2), 120--126. Recuperado de: \url{https://people.csail.mit.edu/rivest/Rsapaper.pdf}

	\hangindent=1.25cm\hangafter=1
	Sánchez, C. A. (2013). \textit{Metodología de la investigación científica y tecnológica} (2.ª ed.). Díaz de Santos.

	\hangindent=1.25cm\hangafter=1
	Schmidt, B., Meier, S., Cremers, C., y Basin, D. (2012). The TAMARIN prover for the symbolic analysis of security protocols. En \textit{Computer Aided Verification (CAV 2012)}, LNCS 7358 (pp. 696--701). Springer. Recuperado de: \url{https://ora.ox.ac.uk/objects/uuid:154668cb-ff09-424a-a492-cf171e21b835}

	\hangindent=1.25cm\hangafter=1
	Shor, P. W. (1997). Polynomial-time algorithms for prime factorization and discrete logarithms on a quantum computer. \textit{SIAM Journal on Computing}, 26(5), 1484--1509. Recuperado de: \url{https://arxiv.org/pdf/quant-ph/9508027.pdf}

	\hangindent=1.25cm\hangafter=1
	Tula, M. I. (2005). \textit{Voto electrónico: Entre la modernización de los sistemas electorales y la desconfianza ciudadana}. Ediciones Dunken.
	
\end{flushleft}